\documentclass[10pt,a4paper,sans]{moderncv}

% Escolha do tema e cor do moderncv
\moderncvstyle{classic}
\moderncvcolor{blue}

% Ajuste das margens da página
\usepackage[scale=0.90]{geometry}
\setlength{\hintscolumnwidth}{2.3cm}

% Informações pessoais
\name{Patrick Serrano}{}
\title{Engenheiro de Software Fullstack Sênior}
\email{patrickcmserrano@gmail.com}
\social[linkedin]{patrickcmserrano}
\social[github]{patrickcmserrano}

\begin{document}
\makecvtitle

\section{Resumo}
Engenheiro de Software Fullstack Sênior com mais de 7 anos de experiência projetando sistemas distribuídos de alta complexidade em ambientes financeiros críticos. Minha formação sólida em Programação Funcional (Clojure, Golang) e Arquitetura Orientada a Eventos me proporciona um modelo mental de engenharia orientado à imutabilidade, consistência de dados e tolerância a falhas --- habilidades que se traduzem diretamente para a construção de microsserviços robustos no ecossistema Node.js/TypeScript. Tenho histórico comprovado em plataformas de pagamento, gateways financeiros e motores de processamento quantitativo, onde alta disponibilidade e resiliência não são opcionais. Integro ativamente ferramentas de IA (Claude, Cursor) ao meu fluxo de desenvolvimento para maximizar qualidade e velocidade de entrega.

\section{Habilidades}
\cvitem{Arquitetura}{Microservices, Event-Driven Architecture, CQRS, Event Sourcing, Clean Architecture, Modular Monolith (Polylith)}
\cvitem{Mensageria}{NATS JetStream, Kafka, RabbitMQ, AWS SQS}
\cvitem{Banco de Dados}{PostgreSQL, Redis (Cache / Sessão), Datomic, DynamoDB, MongoDB}
\cvitem{Testes}{Jest, TDD, Testes de Integração, Testes de Contrato}
\cvitem{Linguagens}{TypeScript, JavaScript, Golang, Clojure, Python}
\cvitem{Frameworks}{Node.js, NestJS, Fastify, React.js, Next.js, Nuxt.js (Vue 3), React Native}
\cvitem{Infra / DevOps}{AWS (SQS, Lambda, ECS), Docker, Kubernetes, CI/CD, Github Actions}
\cvitem{IA no Workflow}{Claude Code (Anthropic), Cursor, Antigravity, Gemini / Google AI Studio}

\section{Experiência Profissional}

\cventry{Jul 2025 -- Atual}{Sistema de Trading Quantitativo (Ark Engine)}{}{}{}{%
Motor de negociação proprietário de altíssima performance focado em baixa latência, tomada de decisão automatizada e resiliência absoluta de dados financeiros.\newline{}
Responsabilidades
\begin{itemize}
    \item Arquitetura e desenvolvimento do zero de um motor de trading de alta frequência, tomando decisões críticas sobre concorrência, consistência transacional e tolerância a falhas em ambiente de produção financeira real;
    \item Migração estratégica de módulos críticos de Python para Golang, alcançando redução drástica de latência no processamento de eventos de mercado em tempo real e aumentando o throughput geral do sistema;
    \item Design e implementação de pipelines de mensageria orientados a eventos com NATS JetStream, garantindo entrega exatamente-uma-vez (exactly-once delivery) e imutabilidade dos registros financeiros, padrão essencial para auditabilidade;
    \item Modelagem de estratégias de cache e gerenciamento de estado em memória para suportar alta taxa de eventos sem degradação de performance sob carga sustentada.
\end{itemize}
Tecnologias: Golang, Python, NATS JetStream, TypeScript, Docker, Event-Driven Architecture
}

\cventry{Dez 2025 -- Atual}{BDL Assetto (Projeto Open-Source / Comunidade)}{}{}{}{%
Sistema de telemetria em tempo real e gerenciamento de estado distribuído para ambiente multiplayer de simulação, com processamento intensivo de dados concorrentes.\newline{}
Responsabilidades
\begin{itemize}
    \item Refatoração de arquitetura monolítica para suportar múltiplos clientes simultâneos, solucionando race conditions críticas e problemas de sincronização de I/O em loops de altíssima frequência;
    \item Implementação de isolamento de estado por cliente (``Per-Pilot View'') com bufferização complexa e janelas de alinhamento de dados, aplicando Design Patterns de concorrência para garantir consistência;
    \item Modelagem de serviços com limiares dinâmicos (Dynamic Thresholds) para tratamento robusto de edge cases em fluxos de dados de sensores em tempo real;
    \item Desenvolvimento de painel de telemetria customizado como camada de observabilidade sobre o sistema nativo, habilitando diagnósticos cirúrgicos em tempo real.
\end{itemize}
Tecnologias: Python, Concorrência, State Management, Telemetria em Tempo Real, Observabilidade
}

\cventry{2022 -- 2025}{Plataforma FinTech (Processamento de Pagamentos em Larga Escala)}{}{}{}{%
Automação de fluxos financeiros críticos e integração de pagamentos com múltiplos gateways e plataformas de e-commerce, operando em alta disponibilidade.\newline{}
Responsabilidades
\begin{itemize}
    \item Desenvolvimento de APIs REST escaláveis e integrações complexas com múltiplos gateways de pagamento (Adyen, Cielo, Getnet) e a plataforma VTEX, garantindo resiliência, rastreabilidade e consistência nas transações;
    \item Refatoração do motor de processamento financeiro de arquitetura acoplada para design modular (Polylith/Microservices), com separação clara de responsabilidades, ampliando cobertura de testes e acelerando ciclos de manutenção;
    \item Implementação de estratégias de cache e filas de processamento assíncrono (via Onyx) para otimizar o throughput em picos de carga, assegurando a integridade do banco de dados Datomic (imutável por design);
    \item Construção de dashboards operacionais com alertas em tempo real sobre a inteligência de roteamento e aprovação de pagamentos, elevando a observabilidade do sistema para times de negócio e engenharia;
    \item Liderança técnica via code reviews rigorosos, pair programming e mentoria de desenvolvedores júnior e pleno, disseminando boas práticas de arquitetura e qualidade de código.
\end{itemize}
Tecnologias: Node.js, Clojure, REST APIs, Arquitetura de Microsserviços, Mensageria, Datomic, AWS
}

\cventry{2020 -- 2021}{Plataforma de E-commerce (CMS em Alta Escala)}{}{}{}{%
Plataforma de gestão de conteúdo para estruturação e exibição de catálogos complexos com alto volume de tráfego.\newline{}
Responsabilidades
\begin{itemize}
    \item Desenvolvimento Fullstack com foco em interfaces reativas e responsivas utilizando React e TypeScript, aplicando Design Patterns de componentização e gerenciamento de estado;
    \item Construção de dashboards reativos para gestão de metadados e ativos de mídia em tempo real com atualizações via WebSocket;
    \item Otimização de infraestrutura AWS com auto-scaling para absorver picos de demanda sem degradação de performance ou disponibilidade.
\end{itemize}
Tecnologias: React, TypeScript, Node.js, ClojureScript, AWS, REST APIs
}

\cventry{2019}{Plataforma Imobiliária}{}{}{}{%
Aplicativo web e mobile para publicidade descentralizada, negociação e venda de imóveis.\newline{}
Responsabilidades
\begin{itemize}
    \item Desenvolvimento fullstack de aplicação mobile e web do zero, com foco em performance de renderização e experiência do usuário;
    \item Modelagem de fluxos de comunicação eficientes entre clientes e corretores de imóveis;
    \item Iteração contínua sobre UX baseada em feedback de usuários reais.
\end{itemize}
Tecnologias: ClojureScript, Fulcro, React Native
}

\cventry{2018}{Projeto FinTech (Gateway de Pagamentos)}{}{}{}{%
Aplicação web para um gateway de pagamentos focada em transações eficientes e seguras.\newline{}
Responsabilidades
\begin{itemize}
    \item Desenvolvimento do front-end da aplicação e lançamento de novas funcionalidades de transação;
    \item Primeiras integrações com sistemas financeiros, consolidando fundamentos de segurança e rastreabilidade de fluxo de pagamento;
    \item Colaboração direta com equipe de produto em ciclos ágeis de entrega.
\end{itemize}
Tecnologias: ClojureScript, Fulcro, REST APIs
}

\section{Formação Acadêmica}
\cventry{2024 -- Atual}{Bacharelado em Engenharia de Software}{Descomplica Faculdade Digital}{}{}{}
\cventry{2015 -- 2018}{Bacharelado em Física}{Universidade Federal Fluminense}{}{}{}

\section{Idiomas}
\cvitemwithcomment{Português}{Nativo}{}
\cvitemwithcomment{Inglês}{Intermediário-Avançado}{}
\cvitemwithcomment{Alemão}{Básico}{}
\cvitemwithcomment{Espanhol}{Básico}{}

\end{document}
